\paragraph{Cyclotomic correction at fixed precision.}
\label{par:thick-integral-correction}
Put
\[
 M=1+\lfloor h/(\ell-1)\rfloor,\qquad
 \sigma_1=\left\lceil 1+(M+\eta)/(\ell-1)\right\rceil .
\]
For a vector \(v\in XN_i(T)\) there is an integral correction
\(\tau(v)\in P_i^{(1)}(T)\) well-defined modulo
\(X^Mq^\eta P_i^{(1)}(T)\), with
\[
  \ell^{\sigma_1}\Phi v-(\gamma-1)\tau(v)
             \in X^Mq^\eta P_i^{(1)}(T).
\]
We use it at an even \(i\), where the odd destination gives
\(q^\eta P_i^{(1)}(T)\subset N_{i+1}(T)\).
There is also a two-step correction with precision \(Xq^{2H}\)
on \(P_i^{(2)}(T)\) and exponent
\(\sigma_*=\lceil1+2H/(\ell-1)^2\rceil\).

Here are the divisions that give these two claims. In either
case solve \((\gamma-1)d=v\) modulo a large power \(X^K\),
over the coefficient field. The degree-by-degree solution
of Paragraph~\ref{par:wach-cut-correction} puts the term in
degree \(y\) in \(\ell^{-e_y}X^yN_i(T)\), where
\[
 e_y=\lceil\ell y/(\ell-1)\rceil.
\]
Take \(d\) with zero
constant coefficient. Write \(J\) for the exponent of \(q\)
in the precision ideal and \(a_y=(\ell-1)y\). In one step
\(\phi(X)^y\) is a unit times \(X^yq^{a_y}\).

If \(a_y<J=\eta\), then \(e_y\le\sigma_1\), since
\(J\le (\ell-1)M\). If \(a_y\ge J\) and \(y<M\),
expand \(q^{a_y-J}=(X+\ell)^{a_y-J}\).
The terms which survive modulo \(X^Mq^J\) are divisible by
\(\ell^{\max(0,\ell y-J-M+1)}\). For \(e_y>\sigma_1\)
we have \(\ell y>J+M+\ell-1\); hence
\[
 \ell y-J-M+1\ge e_y-\sigma_1.
\]

For \(a_y\ge J\) and \(y\ge M\) the entire term vanishes.

For two steps use \(M=1\), \(J=2H\), and
\(\phi^2(X)^y\) associated to
\(X^yq^{a_y}p_1^{a_y}\).
If \(a_y\ge J\) the terms with \(y>0\) vanish. Otherwise
\(p_1\in(\ell,q^\ell)\), by reduction modulo \(\ell\).
After factoring out \(X^yq^{a_y}\), the terms of the binomial
expansion that can survive have \(q^\ell\) used at most
\(\lfloor(J-a_y-1)/\ell\rfloor\) times. If
\(e_y>\sigma_*\), then
\(a_y>J/\ell+(\ell-1)^2/\ell\), and
\[
 a_y-\frac{J-a_y-1}{\ell}
      \ge \frac{\ell a_y-J}{(\ell-1)^2}
      \ge e_y-\sigma_* .
\]
The last inequality uses the defining roundings. Thus the
required division is integral. Here one may take \(K>\max(M,J)\), since
\(\phi^j(X)^K\) then belongs to the precision ideal in both
cases.

The integral coefficient quotient by \(X^Mq^J\) is free
over \(O\) by monic Weierstrass division. For either free Frobenius span \(P\) and its precision submodule \(P_{\rm prec}\), saturation gives
\[
 (P_{\rm prec}[1/\ell])\cap P=P_{\rm prec}.
\]

The resulting representative of \(\ell^{\sigma_1}\Phi d\)
or \(\ell^{\sigma_*}\Phi^2d\) satisfies the claimed
cyclotomic congruence, by commutation and saturation.
Uniqueness of \(d\) modulo \(X^K\) over the coefficient field
follows from the nonzero diagonal \(\chi(\gamma)^{(\ell-1)y}-1\)
in positive degrees. Hence the corrections are linear modulo
the precision and natural in maps of crystalline lattices.
They may be made linear on a free family by choosing
representatives on its basis. In one step the correction to a
vector of \(X^MN_i(T)\) can be taken zero. The calculation of
\(d\) in that case starts at degree \(M\), and
\[
 (\ell-1)M\ge\eta.
\]
For \(v=(\gamma-1)w\) with
\(w\in XN_i(T)\) one can use \(d=w\).

\paragraph{Second restriction and the alternating sum.}
Write \(N_i=N_i(T_2)\) for this construction. For even \(c\)
set
\[
 \xi_c=\ell^{s_1}z_c+(\gamma-1)m_c\in XN_c.
\]

Choose \(t_c=\tau(\xi_c)\) linearly on the restricted family.
At the odd letter \(c+1\) put
\[
 L_c=\ell^{\sigma_1}\Phi m_c+
             \ell^{s_1+\sigma_1}y_{c+1}-t_c,\qquad
 \Xi_c=\ell^{s_1+\sigma_1}z_{c+1}+(\gamma-1)L_c .
\]
The Herr relation gives
\(\Xi_c=\ell^{\sigma_1}\Phi\xi_c-(\gamma-1)t_c\).

Consequently \(\Xi_c\in X^MN_{c+1}\). It lies more precisely
in \(X^Mq^\eta P_c^{(1)}(T_2)\).
To make the next difference integral, write
\[
 U_{c+1}=\ell^R L_c-\ell^{s_1+\sigma_1}\mu_{c+1}
        =\ell^{\sigma_1}\Phi u_c+\ell^{s_1+\sigma_1}A_{c+1}
           -\ell^R t_c .
\]
The signed height and \eqref{eq:thick-extension-data} show that \(q^\eta U_{c+1}\in N_{c+1}\).
Applying \(\gamma-1\) gives \(\ell^R\Xi_c\) minus
\(\ell^{s_1+\sigma_1}(\ell^R z_{c+1}+(\gamma-1)\mu_{c+1})\).
This is in \(N_{c+1}\). Its only possible Laurent obstruction
is again degree zero, this time in
\(\mathrm{Fil}^{-\eta}D_{c+1}/\mathrm{Fil}^0D_{c+1}\).
Only central jumps contribute, so taking one kernel of rank
at most \(m\) per odd destination makes each \(U_{c+1}\) integral.
There are at most \(n\) restrictions including the first ones.

On the resulting family define at each even \(c\)
\[
 w_c=\Phi L_{c-2}+\ell^{s_1+\sigma_1}y_c-\ell^{\sigma_1}m_c,
 \qquad Y_c=\Phi w_c+t_c .
\]
All indices are cyclic, including when the word has length two.
There are three useful tests on \(w_c\). First,
\[
 \ell^R w_c=\Phi U_{c-1}-\ell^{\sigma_1}u_c
                        +\ell^{s_1+\sigma_1}A_c .
\]
The height bounds place the right side, after multiplication
by \(q^{B_c}\), in \(N_c\). Multiplication instead by \(q^\eta\)
places it in \(P_{c-1}^{(1)}(T_2)\) after a possible
additional factor \(\ell^R\), by \eqref{eq:thick-extension-data}.
Both lattices are \(\ell\)-saturated inside the integral
\'etale module.

For the Frobenius span, reduce in a Wach basis
and extend scalars by the finite free Frobenius coefficient map.
The original basis stays independent over the residue Laurent
field because the linearized Frobenius on the \'etale module
is invertible. Cancel these powers of \(\ell\).
Second, \((\gamma-1)w_c=\Phi\Xi_{c-2}-\ell^{\sigma_1}\xi_c\)
belongs to \(XN_c\), since
\(\Phi(X^MN)\subset q X^MN\). Third,
\[
                (\gamma-1)Y_c=\Phi^2\Xi_{c-2}-\Xi_c
                                 \in X^MN_{c+1}.
\]
These calculations give exactly the pole, first pullback
and cyclotomic tests on \(w_c\) used by the even-letter estimates.

Define \(Z_c=(\Phi-1)^{-1}(\gamma-1)Y_c\) in \(X^MN\),
where the inverse exists by contraction on the product of all
letters. Then the sum of \((Y_c,Z_c)\) over even \(c\)
represents \(\ell^{s_1+\sigma_1}(y,z)\). Indeed let
\[
  B=\sum_{c\text{ even}}\bigl(L_c+\Phi L_c
                                  +\ell^{s_1+\sigma_1}y_{c+2}\bigr).
\]
Summing first components gives
\(\sum_cY_c=\ell^{s_1+\sigma_1}y+(\Phi-1)B\)
by invariance of the even labels under the shift by two.
For the second components
\(\ell^{s_1+\sigma_1}z+(\gamma-1)B
 =\sum_c(\Xi_c+\Phi\Xi_c)\in X^MN\).
The Herr relation and uniqueness of the inverse prove the same
coboundary identity for the whole pair.
\label{cor:inert-thick-herr-1}

\paragraph{Integral comparison with the lifted vectors.}
For an admissible \(w\) at an even \(c\), use the first
correction to form
\[
             V(w)=\ell^{\sigma_1}\Phi w
                          -\tau((\gamma-1)w).
\]
Its cyclotomic defect lies in \(X^MN\). The second component is uniquely determined at this precision by
\[
 z=(\Phi-1)^{-1}(\gamma-1)V(w),\qquad z\in X^MN.
\]

Its first component has denominators in finitely many
cyclotomic factors, by the height bounds, and the second
belongs to \(X^MN\). The crystalline criterion proved in
Section~\ref{sec:inert-long-orbit} therefore makes its
integral Herr class finite. A change of representative of
the correction changes both components by a coboundary.
Indeed if the first component of a pair is in \(X^MN\)
along with its second component, the first is the
\((\Phi-1)\)-image of a vector in \(X^MN\). Commutation
and injectivity on \(X^MN\) give the other component of
the coboundary. This also shows that the resulting class
depends linearly on \(w\) and is natural in the marked map.
Writing \([V(w)]\) for the resulting class, the last property of the correction gives
\[
 w\in XN_c\quad\Longrightarrow\quad[V(w)]=0.
\]
Thus integral congruence modulo \(XN_c\) is
sufficient for comparing classes.

For \(w=w_c\), the defect at letter \(c\) is
\(\Phi\Xi_{c-2}-\ell^{\sigma_1}\xi_c\).
Its first summand lies in \(X^MN_c\), on which the
correction is zero. Its second summand has correction
\(-\ell^{\sigma_1}t_c\) modulo the precision.
Thus
\[
 V(w_c)\equiv\ell^{\sigma_1}Y_c\pmod{X^MN}.
\]

The resulting finite classes, denoted by \(\kappa_c\), sum
to \(\ell^{s_1+2\sigma_1}\) times the original class.
The displayed cocycle uses \(\Phi-1\) and \(\gamma-1\)
in this order and has the displayed plus sign. If the
opposite convention for the extension class is used, replace
the sum of the source lifts by its negative.

In particular, if
\[
 \ell^e w_c-f(u_c')\in XN_c(T_2)
\]
for an admissible vector \(u_c'\) in the source, the
finite class attached to \(u_c'\) maps to
\(\ell^e\kappa_c\). All these comparisons take place on
the original rank block and the same map; they do not
require exactness of the Wach functor on the right or a
choice of integral splitting of the surjection.
\label{cor:inert-thick-herr-2}
